% ============================================================
% BAB VI - EVALUASI
% Template (ITB STI TA + report.tex):
%   - VI.1 Metode Evaluasi              : prosedur uji per tujuan T-1..T-4
%   - VI.2 Hasil Evaluasi               : status pemenuhan + metrik
%   - VI.3 Pembahasan Hasil Evaluasi    : interpretasi, keterbatasan, tindak lanjut
% Rantai keterhubungan: T-N <-> VI.1.N <-> VI.2.N <-> VI.3.N <-> Kesimpulan ke-N
% Use case crypto diperkenalkan di sini sebagai instans verifikasi.
% Evaluasi pada lingkup penerimaan fungsional dan struktural node tunggal (B-2).
% Karakterisasi beban kuantitatif multi-node adalah arah lanjutan pada Bab VII.
\cleardoublepage
\chapter{EVALUASI}
\label{chap:evaluasi}

Bab ini memaparkan evaluasi platform terhadap kebutuhan fungsional KF-01 sampai KF-14 dan kebutuhan nonfungsional KNF-01 sampai KNF-12 yang disusun pada Bab~\ref{chap:analisis-masalah}. Evaluasi disusun mengikuti empat tujuan platform T-1 sampai T-4 sehingga setiap hasil dapat dipetakan kembali pada tujuan, rumusan masalah, dan butir kesimpulan dengan indeks yang sama. Pengujian dilakukan pada satu \textit{use case} domain nyata sebagai instans verifikasi, yaitu pasar aset kripto dengan data Coinbase, agar perilaku platform dapat diamati pada beban dan pola data yang realistis tanpa mengubah sifat \textit{domain-agnostic} platform. Penyajian dibagi menjadi metode evaluasi, hasil evaluasi, dan pembahasan, dengan lingkup penerimaan fungsional dan struktural pada lingkungan node tunggal sesuai batasan B-2, sedangkan karakterisasi beban kuantitatif pada klaster multi-node menjadi arah pengembangan yang dirinci pada Bab~\ref{chap:penutup}.

\section{Metode Evaluasi}
\label{sec:metode-evaluasi}

Metode evaluasi mengikuti tiga dimensi yang melekat pada empat tujuan platform, yaitu pemenuhan fungsional, pemenuhan nonfungsional, dan tata kelola. Pemenuhan fungsional menilai apakah setiap KF terbukti melalui skenario uji yang dapat dijalankan ulang, pemenuhan nonfungsional menilai apakah setiap KNF mencapai target metrik yang ditetapkan, dan tata kelola menilai apakah \textit{lineage} tertelusur, kontrak data ditegakkan, serta kebijakan akses dapat diaudit. Setiap skenario diberi penanda SK-F untuk kebutuhan fungsional dan SK-N untuk kebutuhan nonfungsional, lalu dikelompokkan di bawah tujuan T-1 sampai T-4 yang relevan. Penomoran tersebut menjaga keterhubungan satu lawan satu antara skenario, kebutuhan, dan tujuan sehingga kegagalan satu skenario dapat segera dirunut ke komponen yang bertanggung jawab.

Skenario dijalankan secara manual pada tahap evaluasi dan tidak menjadi bagian dari \texttt{make phase-full} maupun rekonsiliasi Argo CD, agar pengukuran tidak mengganggu jalur produksi platform. Pembangkit beban dijalankan dari \textit{namespace} \texttt{platform-test} yang terpisah, sedangkan instrumen pengukuran memanfaatkan komponen observabilitas yang sudah terpasang. Instrumen utama mencakup k6 untuk beban \textit{HTTP} dan \textit{gRPC}, Chaos Mesh untuk injeksi \textit{fault}, Great Expectations untuk kontrak kualitas, serta Prometheus, Grafana, Grafana Tempo, dan OpenCost untuk metrik, \textit{trace}, dan biaya. Berkas skenario disimpan pada direktori \texttt{experiments/} yang sejajar dengan \texttt{platform/} dan \texttt{use-case-crypto/} sebagai lapis verifikasi terpisah, dengan subdirektori \texttt{load/} untuk beban latensi dan \textit{throughput}, \texttt{chaos/} untuk injeksi \textit{fault}, \texttt{drift/} untuk pengujian deteksi \textit{drift}, \texttt{feast/} untuk penyajian fitur \textit{online}, \texttt{parity/} untuk paritas \textit{batch} dan \textit{stream}, \texttt{lineage/} untuk penelusuran \textit{lineage}, serta \texttt{etl/} untuk pemeriksaan pendaratan data. Susunan ini membuat prosedur uji dapat ditelusuri ulang bersama kode tanpa mengubah sistem yang diuji.

\subsection{Instans Verifikasi dan Lingkungan Uji}
\label{subsec:instans-verifikasi}

Instans verifikasi adalah \textit{use case} pasar aset kripto dengan data Coinbase yang dialirkan melalui dua jalur ingestasi yang saling melengkapi. Jalur \textit{real-time} memakai layanan \texttt{websocket-collector} yang berlangganan kanal \textit{ticker} dan \textit{matches} Coinbase pada resolusi per detik hingga per transaksi, lalu \texttt{stream-processor} berbasis Flink mengagregasinya menjadi fitur berjendela waktu pada tabel \texttt{bronze.crypto\_ohlcv\_features}. Jalur historis memakai layanan \texttt{rest-collector} yang menarik \textit{candle} OHLCV Coinbase pada granularitas satu menit sebagai garis dasar sekaligus pemulihan celah historis sejak 1 Maret 2026 mengikuti anggaran penyimpanan pada lingkungan node tunggal sesuai batasan B-2, lalu mendaratkannya pada tabel \texttt{bronze.crypto\_ohlcv}; \textit{endpoint} \textit{candle} REST Coinbase tidak menyediakan granularitas di bawah satu menit sehingga resolusi per detik hanya berlaku pada jalur \textit{real-time}. Pemilihan domain ini didasari karakteristiknya yang menuntut \textit{throughput} ingestasi tinggi, latensi penyajian rendah, dan distribusi data nonstasioner, sehingga keempat \textit{sub-sistem} platform tertekan pada kondisi yang menantang. Seluruh kode dan konfigurasi spesifik domain berada pada direktori \texttt{use-case-crypto/} yang terpisah dari \texttt{platform/}, mencakup layanan \texttt{websocket-collector} dan \texttt{rest-collector} sebagai produsen \textit{event}, \texttt{stream-processor} berbasis Flink, lapis \textit{batch} berbasis dbt, definisi Feast, serta jembatan \textit{inference} dan dasbor. Pemisahan ini menegaskan bahwa platform tetap \textit{domain-agnostic} dan \textit{use case} hanya berperan sebagai muatan uji yang dapat diganti tanpa menyentuh lapis kendali.

Konfigurasi spesifik domain diterapkan melalui \textit{overlay} Kustomize yang menambahkan prefiks \texttt{crypto-} pada sumber daya basis platform, sehingga \textit{InferenceService} basis bernama \texttt{predictor} menjadi \texttt{crypto-predictor} ketika di-\textit{deploy}. \textit{Overlay} tersebut hanya menetapkan prefiks nama, pemetaan registri citra, jumlah \textit{replica} awal, serta batas \textit{autoscaling} untuk lingkungan node tunggal, sedangkan parameter sumber data berada pada ConfigMap basis sehingga definisi dasar tidak perlu disentuh, sebagaimana ditunjukkan pada Listing~\ref{lst:overlay-usecase}. Sesuai invarian node tunggal, setiap \textit{workload} berjalan dengan satu \textit{replica} saat \textit{idle} dan diskalakan oleh HPA serta KEDA hanya ketika beban nyata muncul, dengan \texttt{maxReplicas} dibatasi tiga agar muat pada satu node. Dengan demikian, penggantian \textit{use case} cukup dilakukan dengan menyiapkan \textit{overlay} baru tanpa mengubah manifes platform, sekaligus menjadi bukti konkret atas pemenuhan KF-12 tentang konfigurasi deklaratif dan KNF-11 tentang portabilitas.

\begin{lstlisting}[caption={Kutipan \textit{overlay} Kustomize \textit{use case}: prefiks \texttt{crypto-}, pemetaan registri citra, \textit{replica} \textit{idle} satu, dan batas HPA node tunggal},label={lst:overlay-usecase},basicstyle=\ttfamily\footnotesize,keywordstyle=\ttfamily,breaklines=true,captionpos=t]
# use-case-crypto/manifests/overlays/local/kustomization.yaml
namePrefix: crypto-
resources:
  - ../../base
  - ../../cross-namespace
images:
  - name: rest-collector
    newName: localhost:5000/crypto-rest-collector
    newTag: latest
  - name: websocket-collector
    newName: localhost:5000/crypto-websocket-collector
    newTag: latest
  # ... validator, analyzer, flink-job, trainer, gateway, dbt-project, dll.
replicas:
  - name: rest-collector
    count: 1
  - name: websocket-collector
    count: 1
patches:
  - target:
      kind: HorizontalPodAutoscaler
    patch: |-
      - op: replace
        path: /spec/maxReplicas
        value: 3
      - op: replace
        path: /spec/minReplicas
        value: 1
\end{lstlisting}

\subsection{Evaluasi Tujuan T-1: Arsitektur Terintegrasi dan GitOps}
\label{subsec:metode-t1}

Evaluasi T-1 menilai apakah seluruh platform dapat dipasang dari satu sumber Git, dikelola secara deklaratif, dan menyediakan otomasi siklus hidup model setara MLOps Level 2. Skenario SK-F-01 mengubah definisi \textit{feature view} pada Gitea lalu mengamati rekonsiliasi Argo CD untuk membuktikan KF-01, sementara SK-F-12 mengubah jumlah \textit{replica} pada manifes untuk membuktikan konfigurasi deklaratif KF-12 yang konvergen tanpa \textit{drift}. Skenario SK-F-06 mempromosikan model melalui Argo Rollouts dengan strategi kanari dan memicu \textit{rollback} pada model bermasalah untuk membuktikan otomasi penerapan KF-06, sedangkan SK-F-13 menjalankan satu \textit{notebook} yang memanggil SDK Feast, MLflow, dan klien KServe untuk membuktikan antarmuka terpadu KF-13. Skenario SK-F-14 menjenuhkan kuota \textit{ClusterQueue} Kueue untuk membuktikan manajemen sumber daya KF-14, dengan \textit{job} di luar anggaran berstatus \textit{Pending}.

Kematangan MLOps Level 2 dinilai melalui tiga jalur otomasi yang saling terhubung. Jalur pertama adalah \textit{pipeline} pelatihan otomatis pada Kubeflow Pipelines yang dipicu data baru atau ambang \textit{drift}, jalur kedua adalah penyebaran model otomatis melalui KServe dan Argo Rollouts, dan jalur ketiga adalah pemantauan berkelanjutan oleh Argo Workflows yang melaporkan metrik ke Prometheus. Reproduksibilitas penyebaran diuji oleh SK-N-11 yang menerapkan ulang seluruh manifes pada \textit{node} k3s baru dan SK-N-12 yang mengukur cakupan uji serta waktu konvergensi \texttt{make phase-full}. Gabungan skenario ini menegaskan bahwa T-1 tidak hanya menyediakan komponen, melainkan satu bidang kendali yang dapat direkonstruksi dan diaudit.

Atribut kualitas yang bersifat lintas komponen turut dievaluasi pada T-1 karena melekat pada bidang kendali terintegrasi, bukan pada satu \textit{sub-sistem} tunggal. Skenario SK-N-03 menaikkan beban dengan \texttt{load/hpa-keda-rampup.js} untuk membuktikan skalabilitas horizontal KNF-03 melalui HPA dan KEDA \textit{ScaledObject}, sedangkan SK-N-04 menyuntikkan \textit{fault} dengan Chaos Mesh untuk membuktikan keandalan KNF-04 pada batas RTO dan RPO. Skenario SK-N-08 mengganti satu komponen melalui antarmuka CRD dan \textit{overlay} Kustomize untuk membuktikan ekstensibilitas KNF-08, SK-N-09 mengakses seluruh konsol melalui satu sesi identitas dex untuk membuktikan kemudahan penggunaan KNF-09, dan SK-N-10 mengukur kompresi serta atribusi biaya melalui OpenCost untuk membuktikan efisiensi sumber daya KNF-10. Pengelompokan ini menempatkan T-1 sebagai tujuan dengan beban pengujian nonfungsional terbanyak karena arsitektur terintegrasi memikul jaminan kualitas yang berlaku menyeluruh pada seluruh komponen.

\subsection{Evaluasi Tujuan T-2: Tata Kelola Data}
\label{subsec:metode-t2}

Evaluasi T-2 menilai apakah katalog metadata, \textit{lineage}, dan kontrol kualitas berjalan otomatis sepanjang \textit{pipeline} data, fitur, dan model. Skenario SK-F-08 menelusuri \textit{lineage} pada DataHub melalui \texttt{lineage/}\allowbreak\texttt{lineage-walk.sh} dari topik sumber sampai keluaran prediksi untuk membuktikan KF-08, dengan peristiwa \textit{lineage} dikirim melalui OpenLineage dan \textit{ingestion} terjadwal. Skenario SK-F-09 menjalankan dua \textit{run} pelatihan pada MLflow lalu memutar ulang dari \textit{snapshot} untuk membuktikan pelacakan eksperimen KF-09, sedangkan SK-F-10 memindahkan model melalui status \textit{staging}, \textit{production}, dan \textit{archived} untuk membuktikan registri dan \textit{audit trail} KF-10. Skenario SK-N-05 menjalankan \textit{checkpoint} Great Expectations dan menyuntikkan pelanggaran skema, rentang nilai, serta kelengkapan untuk membuktikan penegakan kontrak kualitas data KNF-05 sepanjang \textit{pipeline}.

Dimensi keamanan dan observabilitas tata kelola dievaluasi melalui dua skenario tambahan. Skenario SK-N-06 menelusuri satu \textit{trace} OpenTelemetry pada Grafana Tempo dari ingestasi sampai prediksi untuk membuktikan observabilitas KNF-06, sehingga metrik, log, dan \textit{trace} dapat dirujuk silang pada satu antarmuka. Skenario SK-N-07 memindai \textit{cluster} dengan Trivy dan memeriksa konfigurasi TLS untuk membuktikan keamanan KNF-07, mencakup mTLS Istio, enkripsi \textit{at rest} MinIO melalui KES dan OpenBao, serta kebijakan admisi Kyverno. Ketiga aspek tersebut menjadikan tata kelola bukan sekadar katalog pasif, melainkan kontrol aktif yang menegakkan kontrak dan kebijakan lintas komponen.

\subsection{Evaluasi Tujuan T-3: Deteksi \textit{Drift} dan Pelatihan Ulang}
\label{subsec:metode-t3}

Evaluasi T-3 menilai apakah deteksi \textit{drift} terotomasi memicu pelatihan ulang sesuai kebijakan dan apakah penyajian fitur menjamin \textit{point-in-time correctness}. Skenario SK-F-07 menerapkan pergeseran distribusi terkendali pada aliran data nyata melalui \texttt{inject\_drift.py}, yang mengonsumsi rekaman pasar tervalidasi lalu menggeser nilainya beberapa sigma sebelum dipublikasikan ulang, kemudian mengamati detektor \textit{drift} dan \textit{CronWorkflow} \texttt{retrain-on-drift} untuk membuktikan pemantauan KF-07. Detektor membandingkan distribusi produksi terhadap distribusi pelatihan dengan \textit{Population Stability Index} dan uji Kolmogorov-Smirnov, lalu memicu \textit{pipeline} pelatihan ketika ambang terlampaui. Skenario SK-F-05 menjalankan \textit{pipeline} pelatihan pada Kubeflow Pipelines yang membaca fitur dari Feast dan menulis model ke MLflow untuk membuktikan otomasi pelatihan KF-05.

Validitas temporal diuji secara khusus karena menjadi pembeda utama dari pendekatan tradisional. Skenario SK-F-03 mengaudit \texttt{get\_historical\_features} pada \textit{dataset} pelatihan untuk membuktikan KF-03, dengan memastikan tidak ada baris yang memakai informasi setelah \textit{event timestamp} acuan. Pengujian ini memakai probe \textit{timestamp} masa depan yang menolak setiap baris bermuatan informasi setelah \textit{event timestamp} acuan sehingga jaminan \textit{zero leakage} ditegakkan langsung pada jalur penyusunan fitur pelatihan. Gabungan skenario ini menegaskan bahwa T-3 menutup dua risiko sekaligus, yaitu degradasi model akibat \textit{drift} dan kebocoran temporal akibat penyusunan fitur yang keliru.

\subsection{Evaluasi Tujuan T-4: Layanan Fitur \textit{Dual-Store}}
\label{subsec:metode-t4}

Evaluasi T-4 menilai apakah layanan fitur \textit{dual-store} menyajikan fitur tabular dan fitur agregat dengan kontrak yang konsisten antara pelatihan dan penyajian, serta menyajikan fitur vektor pada indeks terpisah dengan ruang \textit{embedding} yang sama pada kedua fase tersebut. Skenario SK-F-02 memeriksa penyajian fitur \textit{online} non-\textit{null} pada Valkey melalui \texttt{feast/}\allowbreak\texttt{online-serving-check.py} sebagai bukti sisi \textit{online} penyajian ganda KF-02, sementara kesetaraan nilainya terhadap penyimpanan \textit{offline} ClickHouse bersandar pada proses \textit{materialize} Feast yang menyalin nilai dari penyimpanan \textit{offline} tersebut ke penyimpanan \textit{online} Valkey sehingga keduanya memuat nilai yang sama. Skenario SK-F-04 menjalankan beban pencarian vektor pada Qdrant melalui \texttt{load/}\allowbreak\texttt{vector-search-latency.js} untuk membuktikan dukungan \textit{vector embeddings} KF-04. Skenario SK-F-11 membandingkan fitur hasil \textit{batch} dan hasil \textit{stream} pada indikator yang sama melalui \texttt{parity/}\allowbreak\texttt{parity-check.sh} untuk membuktikan pemrosesan ganda KF-11 dengan definisi fitur yang konsisten. Skenario SK-F-02 dan SK-F-11 menelusuri kontrak fitur tabular dari sisi kebenaran nilai, sedangkan SK-F-04 menegaskan keterbentukan jalur penyajian vektor, sebelum beban tinggi diterapkan.

Kinerja layanan fitur diuji oleh kelompok skenario nonfungsional yang khusus mengukur latensi dan \textit{throughput} jalur penyajian. Skenario SK-N-01 mengukur latensi \textit{feature serving} dan \textit{vector search} dengan k6 melalui \texttt{load/}\allowbreak\texttt{feast-online-latency.js} dan SLO Sloth untuk membuktikan KNF-01, dengan sasaran latensi p99 pada orde sub-detik. Skenario SK-N-02 mengukur \textit{throughput serving} dan \textit{stream} dengan k6 melalui \texttt{load/}\allowbreak\texttt{feast-throughput.js} serta \texttt{kafka-producer-perf-test} untuk membuktikan KNF-02, sambil memantau \textit{lag} konsumer pada beban puncak. Kedua skenario tersebut menempatkan lapis fitur sebagai titik temu seluruh jalur data sehingga kebenaran nilai fitur menjadi prasyarat bagi kualitas layanan yang diukur, sedangkan atribut kualitas lintas komponen ditangani pada T-1.

\subsection{Matriks Cakupan dan Uji Penerimaan}
\label{subsec:cakupan}

Pemetaan antara skenario dan kebutuhan bersifat banyak-ke-banyak, sebab satu skenario kerap menegaskan lebih dari satu kebutuhan dan satu kebutuhan dapat diperkuat oleh beberapa skenario. Sebagai contoh, SK-F-04 yang menguji latensi pencarian vektor sekaligus menegaskan KNF-01, sedangkan SK-N-01 yang mengukur latensi p99 penyajian turut memperkuat KF-04 pada sisi kualitas layanan vektor. Tabel~\ref{tab:uji_penerimaan} merangkum prosedur uji dan kriteria penerimaan untuk setiap skenario yang dikelompokkan di bawah tujuan T-1 sampai T-4. Susunan tersebut memastikan seluruh empat belas kebutuhan fungsional dan dua belas kebutuhan nonfungsional memiliki sekurang-kurangnya satu skenario uji yang konkret. Katalog perintah yang dapat dijalankan ulang untuk setiap skenario, termasuk preambel pemuatan konfigurasi \texttt{config.crypto.env} dan pembuatan \textit{namespace} \texttt{platform-test}, dirangkum pada Lampiran~\ref{appx:verifikasi} sehingga setiap baris Tabel~\ref{tab:uji_penerimaan} dapat ditelusuri ke perintah konkret yang menghasilkan buktinya.

{
\footnotesize
\setlength{\tabcolsep}{3pt}
\begin{longtable}{|>{\centering\arraybackslash}m{1.4cm}|>{\centering\arraybackslash}m{1.75cm}|m{5.9cm}|m{4cm}|}
\caption{Skenario Uji Penerimaan per Tujuan Platform}
\label{tab:uji_penerimaan} \\
\hline
\textbf{Skenario} & \textbf{Kebutuhan} & \multicolumn{1}{>{\centering\arraybackslash}m{5.9cm}|}{\textbf{Prosedur Uji}} & \multicolumn{1}{>{\centering\arraybackslash}m{4cm}|}{\textbf{Kriteria Penerimaan}} \\
\hline
\endfirsthead
\hline
\textbf{Skenario} & \textbf{Kebutuhan} & \multicolumn{1}{>{\centering\arraybackslash}m{5.9cm}|}{\textbf{Prosedur Uji}} & \multicolumn{1}{>{\centering\arraybackslash}m{4cm}|}{\textbf{Kriteria Penerimaan}} \\
\hline
\endhead
\endfoot
\hline
\endlastfoot

\multicolumn{4}{|l|}{\textbf{Tujuan T-1: arsitektur platform terintegrasi dan praktik GitOps}} \\
\hline
SK-F-01 & KF-01 & Mengubah definisi \textit{feature view} pada repositori Gitea lalu mengamati rekonsiliasi Argo CD, kemudian memverifikasi \texttt{feast feature-views list}. & \textit{Feature view} baru aktif tanpa intervensi manual dan perubahan terlacak pada riwayat Git. \\
\hline
SK-F-06 & KF-06 & Mempromosikan model melalui Argo Rollouts dengan strategi kanari, lalu menyuntikkan model bermasalah untuk memicu \textit{rollback}. & Model valid dipromosikan dan model bermasalah memicu \textit{rollback} otomatis tanpa kehilangan permintaan. \\
\hline
SK-F-12 & KF-12 & Mengubah jumlah \textit{replica} pada manifes YAML di Gitea, lalu mengamati rekonsiliasi Argo CD. & Status \textit{cluster} konvergen ke definisi Git dan tidak menyisakan \textit{drift}. \\
\hline
SK-F-13 & KF-13 & Menjalankan notebook satu \textit{environment} uv yang memuat SDK Feast, MLflow, dan klien \textit{inference} KServe v2, lalu memanggil ketiganya sebagai operasi baca independen. & Ketiga klien termuat dan berjalan dari satu \textit{environment} uv tanpa konfigurasi tambahan per layanan. \\
\hline
SK-F-14 & KF-14 & Mengirim sejumlah \textit{job} pelatihan melebihi kuota \textit{ClusterQueue} Kueue, lalu mengamati status admisi. & \textit{Job} di luar anggaran kuota berstatus \textit{Pending}. Admisi mengikuti \textit{nominalQuota} dan \textit{borrowingLimit}. \\
\hline
SK-N-03 & KNF-03 & Menaikkan beban dengan \texttt{load/hpa-keda-rampup.js} sambil mengamati HPA, KEDA \textit{ScaledObject}, dan jumlah \textit{pod}. & Jumlah \textit{replica} naik mengikuti beban dan \textit{throughput} bertambah mendekati linear. \\
\hline
SK-N-04 & KNF-04 & Menyuntikkan \textit{fault} dengan Chaos Mesh (\texttt{pod-chaos} dan \texttt{network-chaos}) lalu mengamati pemulihan. & Beban kerja pulih otomatis dengan RTO dan RPO di bawah ambang yang ditetapkan. \\
\hline
SK-N-08 & KNF-08 & Mengganti satu komponen (misalnya \textit{runtime} penyajian) melalui antarmuka CRD dan \textit{overlay} Kustomize. & Penggantian terisolasi pada antarmuka tanpa mengubah komponen lain. \\
\hline
SK-N-09 & KNF-09 & Mengakses seluruh konsol platform melalui satu sesi identitas dex serta memeriksa ketersediaan dokumentasi operasional pada repositori Gitea. & Seluruh konsol dapat dicapai melalui satu jalur autentikasi tanpa login ulang dan dokumentasi prosedur inti tersedia. \\
\hline
SK-N-10 & KNF-10 & Mengukur rasio kompresi ClickHouse, utilisasi sumber daya, dan atribusi biaya \textit{workload} pada OpenCost. & Rasio kompresi dan utilisasi memenuhi target serta biaya per \textit{workload} dapat diatribusikan. \\
\hline
SK-N-11 & KNF-11 & Menerapkan ulang seluruh manifes GitOps pada \textit{node} k3s baru dari satu sumber Git. & Seluruh komponen terekonstruksi tanpa intervensi manual selain \textit{bootstrap} Argo CD. \\
\hline
SK-N-12 & KNF-12 & Menjalankan \texttt{make usecase-crypto-test} untuk cakupan uji dan \texttt{make phase-full} pada \textit{cluster} bersih sambil mengukur waktu konvergensi. & Cakupan uji melewati ambang yang ditetapkan dan \texttt{phase-full} konvergen tanpa langkah manual. \\
\hline

\multicolumn{4}{|l|}{\textbf{Tujuan T-2: sub-sistem tata kelola data}} \\
\hline
SK-F-08 & KF-08 & Menelusuri \textit{lineage} pada DataHub dari topik sumber hingga keluaran prediksi melalui antarmuka katalog. & Rantai \textit{lineage} utuh dari sumber data sampai prediksi tersaji pada satu graf metadata. \\
\hline
SK-F-09 & KF-09 & Menjalankan dua \textit{run} pelatihan lalu mencatatnya pada MLflow, kemudian memutar ulang \textit{run} dari \textit{snapshot} data dan parameter tercatat. & \textit{Run} tercatat lengkap dan pemutaran ulang menghasilkan selisih metrik di bawah satu persen. \\
\hline
SK-F-10 & KF-10 & Memindahkan model melalui status \textit{none}, \textit{staging}, \textit{production}, dan \textit{archived} pada registri MLflow. & Transisi mengikuti aturan yang ditetapkan dan \textit{audit trail} terekam lengkap. \\
\hline
SK-N-05 & KNF-05 & Menjalankan \textit{checkpoint} Great Expectations pada \textit{batch} data, lalu menyuntikkan pelanggaran skema, rentang nilai, dan kelengkapan untuk menguji penegakan kontrak. & Setiap pelanggaran kontrak kualitas terdeteksi dan \textit{batch} cacat ditolak sebelum masuk lapis hilir. \\
\hline
SK-N-06 & KNF-06 & Menelusuri satu \textit{trace} OpenTelemetry pada Grafana Tempo dari ingestasi hingga prediksi. & Satu \textit{trace} menautkan lintasan kritis dan metrik, log, serta \textit{trace} dapat dirujuk silang. \\
\hline
SK-N-07 & KNF-07 & Memindai \textit{cluster} dengan Trivy dan memeriksa konfigurasi TLS pada \textit{endpoint} layanan. & TLS 1.3 \textit{in transit} dan AES-256 \textit{at rest} terverifikasi tanpa kerentanan kritis terbuka. \\
\hline

\multicolumn{4}{|l|}{\textbf{Tujuan T-3: deteksi \textit{drift} dan pelatihan ulang}} \\
\hline
SK-F-03 & KF-03 & Mengaudit \texttt{get\_historical\_features} pada \textit{dataset} pelatihan untuk memastikan jendela \textit{point-in-time} dipatuhi. & Tidak ada baris yang memakai informasi setelah \textit{event timestamp} acuan. \\
\hline
SK-F-05 & KF-05 & Menjalankan \textit{pipeline} pelatihan pada Kubeflow Pipelines yang membaca fitur dari tabel \textit{gold} dan menulis model ke MLflow. & \textit{Pipeline} berjalan dari ujung ke ujung dan model tercatat sebagai artefak \textit{run} pada MLflow yang dapat disajikan KServe. \\
\hline
SK-F-07 & KF-07 & Menyuntikkan \textit{drift} terkendali pada satu fitur data nyata, lalu mengamati detektor \textit{drift} dan \textit{CronWorkflow retrain-on-drift}. & \textit{Drift} terdeteksi saat ambang PSI atau KS terlampaui dan \textit{pipeline} pelatihan ulang terpicu otomatis. \\
\hline

\multicolumn{4}{|l|}{\textbf{Tujuan T-4: layanan fitur \textit{dual-store}}} \\
\hline
SK-F-02 & KF-02 & Memeriksa penyajian fitur \textit{online} pada Valkey mengembalikan nilai non-\textit{null} melalui \texttt{online-serving-check.py}. & Penyajian \textit{online} mengembalikan nilai non-\textit{null} dan kesetaraan nilai \textit{online} dan \textit{offline} terpenuhi secara struktural melalui \textit{materialize}. \\
\hline
SK-F-04 & KF-04 & Menjalankan beban pencarian vektor pada Qdrant melalui \texttt{load/vector-search-latency.js} dengan k6. & Latensi pencarian vektor berada pada orde sub-detik; pengujian \textit{recall} dilakukan saat koleksi vektor terisi. \\
\hline
SK-F-11 & KF-11 & Membandingkan fitur hasil \textit{batch} (dbt) dan hasil \textit{stream} (Flink) untuk indikator yang sama. & Nilai fitur \textit{batch} dan \textit{stream} setara dalam toleransi yang ditetapkan. \\
\hline
SK-N-01 & KNF-01 & Menjalankan \texttt{load/feast-online-latency.js} dan beban pencarian vektor dengan k6 sambil memantau SLO Sloth. & Latensi p99 berada pada orde sub-detik pada \textit{feature serving} dan \textit{vector search}. \\
\hline
SK-N-02 & KNF-02 & Menjalankan \texttt{load/feast-throughput.js} untuk \textit{serving} dan \texttt{kafka-producer-perf-test} untuk \textit{stream}, sambil memantau \textit{lag}. & \textit{Throughput} mencapai target dan \textit{lag} konsumer tetap terkendali pada beban puncak. \\
\hline

\end{longtable}
}


\section{Hasil Evaluasi}
\label{sec:hasil-evaluasi}

Hasil evaluasi dilaporkan sebagai uji penerimaan fungsional dan struktural pada lingkungan node tunggal sesuai batasan B-2. Status setiap kebutuhan dinyatakan dalam tiga tingkat, yaitu terverifikasi struktural ketika jalur kendali terbentuk dan terekonsiliasi sesuai rancangan, terverifikasi fungsional ketika skenario uji dijalankan dan keluarannya memenuhi kriteria penerimaan, serta terinstrumentasi ketika instrumen ukur telah terpasang dan jalur melayani permintaan sementara profil beban kuantitatif penuh berada pada lingkup pengembangan lanjutan. Pembedaan tiga tingkat ini menjaga agar setiap klaim sepadan dengan bukti yang tersedia dan tidak melampaui hasil yang benar-benar diamati. Tabel~\ref{tab:evaluasi_target} merangkum target, instrumen ukur, dan status setiap kebutuhan, sedangkan karakterisasi beban kuantitatif pada klaster multi-node, seperti distribusi latensi p99 dan \textit{throughput} puncak, dirinci sebagai arah pengembangan pada Bab~\ref{chap:penutup}.

Pada T-1, verifikasi struktural menunjukkan seluruh komponen berhasil dipasang dan direkonsiliasi oleh Argo CD dari satu repositori Gitea, perubahan \textit{replica} dan definisi fitur konvergen tanpa \textit{drift}, dan rekonstruksi pada \textit{node} k3s baru berhasil tanpa intervensi manual selain \textit{bootstrap} awal. Pada T-2, peristiwa \textit{lineage} dari emitter OpenLineage Airflow dan Spark serta konektor \textit{ingestion} DataHub untuk Feast, MLflow, ClickHouse, dan Kafka muncul pada satu graf metadata dengan kosa kata seragam, transisi status model tercatat pada MLflow, dan kontrol keamanan mTLS serta enkripsi \textit{at rest} aktif sesuai konfigurasi. Pada T-3, \textit{control loop} pelatihan ulang terpicu otomatis ketika injeksi \textit{drift} terkendali diterapkan pada satu fitur, sehingga rantai deteksi hingga pemicuan pelatihan ulang terbukti berfungsi pada jalur kendali. Pada T-4, jalur penyajian fitur \textit{online}, \textit{offline}, dan vektor berhasil dibentuk dan melayani permintaan, sementara karakterisasi latensi p99 dan \textit{throughput} pada beban tinggi dirinci sebagai arah pengembangan pada klaster multi-node di Bab~\ref{chap:penutup}.

Beberapa skenario pada instans verifikasi menghasilkan keluaran kuantitatif yang dapat dijalankan ulang melalui katalog perintah pada Lampiran~\ref{appx:verifikasi}, sebelum status ringkas setiap kebutuhan dirangkum pada Tabel~\ref{tab:evaluasi_target}. Keluaran tersebut beserta konteks dan batasan penafsirannya dirangkum sebagai berikut.

\vspace{0.5em}
\begin{enumerate}
    \item \textbf{Pemicuan pelatihan ulang oleh deteksi \textit{drift} (KF-07)}

    Injeksi pergeseran terkendali sebesar dua sigma pada fitur \texttt{return\_24h} menghasilkan \textit{Population Stability Index} sebesar 18,107 yang jauh melampaui ambang 0,2 dan statistik Kolmogorov-Smirnov sebesar 0,815 di atas ambang 0,15, sehingga detektor menandai \textit{drift} dan memicu satu eksekusi \textit{CronWorkflow} \texttt{retrain-on-drift}. Angka ini berasal dari injeksi terkendali pada aliran tervalidasi melalui \texttt{experiments/drift/inject\_drift.py --sigma 2.0}, sehingga membuktikan keterpicuan rantai pelatihan ulang dan bukan ketepatan deteksi terhadap pergeseran distribusi alami berdurasi panjang.

    \vspace{0.5em}

    \item \textbf{Pendaratan data nyata pada lapis medali}

    Pemeriksaan pendaratan melaporkan sekitar 307.556 baris pada \texttt{bronze.crypto\_ohlcv} dan \texttt{silver.stg\_ohlcv} serta sekitar 307.522 baris pada \texttt{gold.fct\_ohlcv\_features} dan \texttt{gold.fct\_training\_data}, mencakup dua pasangan instrumen pada jendela 1 Maret sampai 16 Juni 2026 dengan resolusi satu menit. Substrat ini menjadi dasar skenario yang membaca lapis \textit{gold} seperti penyusunan fitur historis (KF-03) dan otomasi pelatihan (KF-05), dijalankan melalui \texttt{experiments/etl/data-landed-check.sh} pada jalur \textit{batch} REST yang masih terus bertambah; pemeriksaan ini menilai pendaratan \textit{batch} dan bukan paritas fitur \textit{batch} dan \textit{stream} (KF-11).

    \vspace{0.5em}

    \item \textbf{Keutuhan \textit{lineage} (KF-08)}

    Penelusuran graf metadata DataHub menemukan 270 \textit{dataset}, 18 \textit{DataJob}, dan 3 \textit{DataFlow} yang saling terhubung dari topik sumber sampai keluaran prediksi, dijalankan melalui \texttt{experiments/lineage/lineage-walk.sh discover}. Cakupan graf ditentukan oleh \textit{namespace} OpenLineage yang dipakai \textit{emitter}, sehingga angka ini mengukur keterhubungan \textit{lineage} yang teremit dan terkatalog, bukan kelengkapan seluruh aset lintas \textit{namespace}.

    \vspace{0.5em}

    \item \textbf{Penyajian pencarian vektor (KF-04)}

    Beban pencarian vektor pada Qdrant mencatat latensi p50 sebesar 4,25 ms dan p95 sebesar 5,02 ms pada 100 iterasi, dijalankan melalui \texttt{k6 run experiments/load/vector-search-latency.js}. Pengukuran ini menegaskan jalur pencarian vektor terpasang dan melayani permintaan pada orde milidetik, bukan kualitas \textit{recall} karena koleksi uji masih kosong, sedangkan nilai p99 sebesar 157 ms merupakan \textit{artifact} fase pembongkaran beban dan bukan latensi tunak.
\end{enumerate}

{
\footnotesize
\setlength{\tabcolsep}{3pt}
\begin{longtable}{|>{\centering\arraybackslash}m{1.7cm}|m{5cm}|m{3.8cm}|m{2.5cm}|}
\caption{Target Metrik, Instrumen Ukur, dan Status Evaluasi per Kebutuhan}
\label{tab:evaluasi_target} \\
\hline
\textbf{Kebutuhan} & \multicolumn{1}{>{\centering\arraybackslash}m{5cm}|}{\textbf{Target atau Kriteria}} & \multicolumn{1}{>{\centering\arraybackslash}m{3.8cm}|}{\textbf{Instrumen Ukur}} & \multicolumn{1}{>{\centering\arraybackslash}m{2.5cm}|}{\textbf{Status}} \\
\hline
\endfirsthead
\hline
\textbf{Kebutuhan} & \multicolumn{1}{>{\centering\arraybackslash}m{5cm}|}{\textbf{Target atau Kriteria}} & \multicolumn{1}{>{\centering\arraybackslash}m{3.8cm}|}{\textbf{Instrumen Ukur}} & \multicolumn{1}{>{\centering\arraybackslash}m{2.5cm}|}{\textbf{Status}} \\
\hline
\endhead
\endfoot
\hline
\endlastfoot

\multicolumn{4}{|l|}{\textbf{Kebutuhan Fungsional}} \\
\hline
KF-01 & \textit{Feature view} terversion dan rekonsiliasi otomatis & Feast \textit{registry}, Argo CD, Gitea & Terverifikasi struktural \\
\hline
KF-02 & Konsistensi nilai \textit{online} dan \textit{offline} & Feast SDK, k6 & Menunggu pengukuran beban \\
\hline
KF-03 & \textit{Zero leakage} pada \textit{point-in-time join} & Feast PIT, probe \textit{leakage} & Menunggu pengukuran beban \\
\hline
KF-04 & \textit{Recall} memenuhi target, latensi sub-detik & Qdrant, k6 & Menunggu pengukuran beban \\
\hline
KF-05 & \textit{Pipeline end-to-end}, model terdaftar & Kubeflow Pipelines, MLflow & Terverifikasi struktural \\
\hline
KF-06 & Promosi kanari dan \textit{rollback} otomatis & KServe, Argo Rollouts, Prometheus & Terverifikasi struktural \\
\hline
KF-07 & \textit{Drift} memicu pelatihan ulang & \textit{drift-detector}, Argo CronWorkflow & Struktural, ketepatan menyusul \\
\hline
KF-08 & \textit{Lineage} data ke fitur ke model utuh & DataHub, OpenLineage & Terverifikasi struktural \\
\hline
KF-09 & \textit{Run} reproducible, selisih ulang kecil & MLflow & Menunggu pengukuran beban \\
\hline
KF-10 & Transisi status dan \textit{audit trail} & MLflow Model Registry & Terverifikasi struktural \\
\hline
KF-11 & Paritas fitur \textit{batch} dan \textit{stream} & dbt, Flink & Menunggu pengukuran beban \\
\hline
KF-12 & Konvergensi deklaratif tanpa \textit{drift} & Kustomize, Argo CD, Gitea & Terverifikasi struktural \\
\hline
KF-13 & Satu sesi memanggil seluruh SDK & SDK Feast, MLflow, KServe & Terverifikasi struktural \\
\hline
KF-14 & \textit{Job} luar kuota \textit{Pending}, admisi sesuai & Kueue \textit{ClusterQueue}, HPA, VPA, KEDA & Terverifikasi struktural \\
\hline

\multicolumn{4}{|l|}{\textbf{Kebutuhan Nonfungsional}} \\
\hline
KNF-01 & Latensi p99 sub-detik \textit{serving} dan vektor & k6, SLO Sloth & Menunggu pengukuran beban \\
\hline
KNF-02 & \textit{Throughput serving} dan \textit{stream} tinggi & k6, \texttt{kafka-producer-}\allowbreak\texttt{perf-test}, KEDA & Menunggu pengukuran beban \\
\hline
KNF-03 & \textit{Replica} dan \textit{throughput} naik mendekati linear & HPA, KEDA, k6 & Menunggu pengukuran beban \\
\hline
KNF-04 & Pemulihan otomatis, RTO dan RPO terjaga & Chaos Mesh, Velero, Sloth & Struktural, RTO/RPO menyusul \\
\hline
KNF-05 & \textit{Zero} kebocoran data temporal & Great Expectations, probe & Struktural, kebenaran data menyusul \\
\hline
KNF-06 & Metrik, log, dan \textit{trace} silang-rujuk & Prometheus, Loki, Tempo, OTel & Terverifikasi struktural \\
\hline
KNF-07 & TLS 1.3, AES-256, dan \textit{audit log} aktif & Istio mTLS, KES, OpenBao, Trivy & Terverifikasi struktural \\
\hline
KNF-08 & Penggantian komponen terisolasi & CRD, Helm, Kustomize & Terverifikasi struktural \\
\hline
KNF-09 & SSO satu sesi tanpa login ulang, dokumentasi tersedia & dex, oauth2-proxy, Gitea & Terverifikasi struktural \\
\hline
KNF-10 & Kompresi, utilisasi, dan atribusi biaya & OpenCost, ClickHouse & Menunggu pengukuran beban \\
\hline
KNF-11 & Rekonstruksi penuh dari Git & Argo CD, Kustomize & Terverifikasi struktural \\
\hline
KNF-12 & Cakupan uji dan \textit{phase-full} konvergen & Tekton CI, Argo CD & Terverifikasi struktural \\
\hline

\end{longtable}
}


\section{Pembahasan Hasil Evaluasi}
\label{sec:pembahasan-evaluasi}

Pembahasan mengaitkan setiap hasil dengan tujuan yang dijawab dan dengan rumusan masalah yang menjadi titik berangkat. Verifikasi struktural pada T-1 menjawab RM-1 dengan menunjukkan bahwa fragmentasi tumpukan teknologi dapat diganti oleh satu bidang kendali GitOps yang dapat direkonstruksi. Verifikasi pada T-2 menjawab RM-2 dengan menunjukkan tata kelola berjalan sebagai layanan bersama, bukan pekerjaan pasca-fakta, sehingga \textit{lineage} dan kontrak kualitas tertanam pada \textit{pipeline}. Verifikasi pada T-3 menjawab RM-3 dengan menunjukkan bahwa deteksi \textit{drift} memicu pelatihan ulang secara otomatis pada jalur kendali, sehingga jawaban atas RM-3 tegak pada tingkat fungsional dan struktural. Verifikasi pada T-4 menjawab RM-4 dengan menunjukkan bahwa layanan fitur \textit{dual-store} terbentuk dan melayani permintaan dengan kontrak fitur yang konsisten, sehingga jawaban atas RM-4 tegak pada tingkat struktural dengan jalur penyajian yang telah terinstrumentasi, sedangkan karakterisasi beban kuantitatifnya menjadi arah pengembangan.

Cakupan evaluasi telah memenuhi syarat keterlacakan, karena seluruh empat belas kebutuhan fungsional dan dua belas kebutuhan nonfungsional memiliki sekurang-kurangnya satu skenario uji pada Tabel~\ref{tab:uji_penerimaan} dan satu baris status pada Tabel~\ref{tab:evaluasi_target}. Pemetaan banyak-ke-banyak memperkuat keyakinan karena beberapa kebutuhan kritis, seperti KNF-01 dan KF-04, ditegaskan oleh lebih dari satu skenario. Dengan demikian, kegagalan pada satu skenario tidak langsung membatalkan klaim pemenuhan kebutuhan terkait, melainkan menyisakan jalur verifikasi lain yang masih dapat ditelusuri. Susunan ini sekaligus menyiapkan struktur kesimpulan pada Bab~\ref{chap:penutup} yang menjawab tepat satu tujuan untuk setiap butir.

Evaluasi ini memiliki tiga keterbatasan yang perlu dicatat secara terbuka.

\vspace{0.5em}
\begin{enumerate}
    \item \textbf{Lingkungan pengujian node tunggal}

    Lingkungan pengujian adalah k3s node tunggal sesuai batasan B-2, sehingga klaim skalabilitas horizontal diukur melalui \textit{multi-replica} dan \textit{autoscaling} pada satu \textit{node}, bukan pada klaster multi-node.

    \vspace{0.5em}

    \item \textbf{Verifikasi \textit{drift} melalui injeksi terkendali}

    Deteksi \textit{drift} diverifikasi melalui injeksi \textit{drift} terkendali pada data nyata sehingga ketepatan terhadap pergeseran distribusi alami berdurasi panjang masih perlu diamati lebih lama, sedangkan karakterisasi beban penuh seperti latensi p99 dan \textit{throughput} puncak menjadi arah pengembangan pada lingkungan multi-node.

    \vspace{0.5em}

    \item \textbf{Verifikasi kesetaraan nilai fitur \textit{dual-store}}

    Kesetaraan nilai fitur antara penyimpanan \textit{online} Valkey dan \textit{offline} ClickHouse pada KF-02 bersifat \textit{structural-by-construction} karena proses \textit{materialize} Feast menyalin nilai dari penyimpanan \textit{offline} ke penyimpanan \textit{online} sehingga keduanya memuat nilai yang sama, sedangkan uji pembandingan nilai eksplisit per entitas melalui SDK Feast masih berupa prosedur manual yang belum terotomasi pada jalur evaluasi, sehingga KF-02 dilaporkan pada tingkat struktural dan penegasan fungsionalnya menjadi tindak lanjut.
\end{enumerate}

Ketiga keterbatasan tersebut menjadi dasar saran tindak lanjut pada Bab~\ref{chap:penutup}, terutama untuk pengujian beban pada lingkungan multi-node dan pengamatan \textit{drift} pada data nyata berdurasi panjang.
